The estimated mismatch of approximately four in ten evaluable employed STEM records in 2025 indicates that employment alone does not ensure occupational alignment. This distinction agrees with the conceptual separation of horizontal mismatch, overqualification and skills mismatch discussed by McGuinness et al. and Somers et al. \cite{McGuinness2018,Somers2019}. Robst and Verhaest et al. further distinguish degree-field and qualification-level mismatch; Freeman and Hirsch examine the knowledge content of jobs \cite{Robst2008,Verhaest2017,Freeman2008}. The result is not directly comparable with the OECD figures introduced earlier because the populations, periods and measurement approaches differ. In particular, the rule adapted from Montt \cite{Montt2017} identifies normative compatibility rather than the actual use of skills in a particular job. Task-based or self-reported measurements \cite{Reis2018,Salas-Velasco2021} could classify some records differently.

The differences across STEM fields and sex support the decision to analyse these groups separately. They are consistent with the relevance of heterogeneity discussed in the literature \cite{Kiss2024,Deming2020}, but the observed differences do not identify their mechanisms. Research on STEM career exits and family trajectories, and UNESCO's work on gender barriers in STEM education, offer complementary contexts without establishing the mechanisms in Ecuador \cite{Hunt2016,Cech2019,UNESCO2017}. The income and employment-condition comparisons also provide context for previous studies of mismatch and labour-market outcomes \cite{Montt2017,Takeuchi2023,BotelloPealoza2016,DeLaTorreB2025,Nordin2010}. Bassi et al. provide a wider regional perspective on education and employment \cite{Bassi2012}. The lower median income and lower adequate-employment share of mismatched records should not be interpreted as causal wage penalties. Differences in age, sector, work hours, employment arrangements and missing income could contribute to these patterns.

The predictive comparison answers the research question without assuming that more complex models are better. The highest ROC-AUC was 0.6522, and AP improved over the class prevalence of 0.4028, but the threshold-dependent results show limited individual classification. Logistic regression provides a competitive benchmark, in agreement with the rationale for empirical comparisons in tabular learning \cite{James2023,Grinsztajn2022,McElfresh2023}. The paired intervals do not establish superiority between the fitted classifiers. Sex, field and province were consistently informative; field is partly involved in defining the outcome, so the decrease after excluding it must be understood in light of the label construction. Importance is a measure of predictive contribution rather than an explanation of causality \cite{Molnar2025}.

\subsection{Limitations}
The main limitation is the broad normative correspondence, which combines fields 05 and 06 and uses three-digit occupation groups; it has not been independently validated for Ecuador. A teacher or manager classified as mismatched may still use STEM knowledge. The absence of an alternative documented rule prevents assessing uncertainty in the label itself. Cross-year household identity is not verified, even though the candidate purge provided a conservative sensitivity check. The 2025 data were inspected during the first analysis and the corrected AP evaluation is not a wholly new holdout. Bootstrap intervals are approximate under the survey design and singleton strata; they do not account for label uncertainty. Income comparisons exclude missing/non-finite values and the nonresponse sentinel 999999, but retain recorded negative incomes (22 in 2025); the interpretation of these values requires confirmation. Nominal values are not adjusted for inflation. The 2021 precision issue also limits exact-degree analysis. Finally, the analysis concerns employed and evaluable graduates only: it cannot measure vacancies, shortages, unemployment risk, or future career trajectories; the OECD distinguishes mismatch from skills shortages and surpluses \cite{OECD2017}. Further research could validate the correspondence with occupational experts or task-level information and assess model performance on a genuinely new annual sample.
